\uuid{CF1k}
\exo7id{4761}
\titre{exo7 4761}
\auteur{quercia}
\organisation{exo7}
\datecreate{2010-03-16}
\isIndication{false}
\isCorrection{false}
\chapitre{Topologie}
\sousChapitre{Topologie des espaces vectoriels normés}
\cpge{mpsi-mp2i.nombres-reels-suites.c, pcsi.nombres-reels-suites.b, ptsi.nombres-reels-suites.b, bcpst1.suites-reelles, mp-mpi.topologie-espaces-normes.a, pc.espaces-vectoriels-normes.a, psi.espaces-vectoriels-normes.a, mp-mpi.topologie-espaces-normes.c, pc.espaces-vectoriels-normes.c, psi.espaces-vectoriels-normes.c}
\module{Analyse}
\niveau{L2}
\difficulte{}

\contenu{
\question{
Soit $(\lambda_n)$ une suite de r{\'e}els strictement positifs.
On lui associe la norme sur $\R[x]$ :
$N(\sum\limits_i a_ix^i) = \sum\limits_i \lambda_i|a_i|$.

Soient $(\lambda_n)$ et $(\lambda'_n)$ deux suites et $N$, $N'$ les normes
associ{\'e}es.
Montrer que $N$ et $N'$ sont {\'e}quivalentes si et seulement si les suites
$\Bigl(\frac{\lambda_n}{\lambda'_n}\Bigr)$ et
$\Bigl(\frac{\lambda'_n}{\lambda_n}\Bigr)$ sont born{\'e}es.
}
}
